% Dirichlet Round 08 English TeX
% Translation of Werke II, pp. 105--114.
\documentclass[11pt]{article}
\usepackage[a4paper,margin=1in]{geometry}
\usepackage[T1]{fontenc}
\usepackage[utf8]{inputenc}
\usepackage{lmodern}
\usepackage{amsmath,amssymb}
\usepackage{microtype}
\setlength{\parindent}{1.2em}
\setlength{\parskip}{0.25em}
\begin{document}

\begin{center}
{\Large G. Lejeune Dirichlet's Works. II.}\par\vspace{1.2em}
{\Large X.}\par\vspace{0.5em}
{\LARGE \textbf{On the Composition of Binary Forms}}\par
{\LARGE \textbf{of the Second Degree.}}\par\vspace{1em}
{\large P. G. Lejeune Dirichlet.}\par\vspace{0.5em}
{\small Commentary, Berlin, Academic Press, 1851.}
\end{center}

\begin{center}
{\large \textbf{On the Composition of Binary Forms of the Second Degree\footnote{In the printing on pp. 155--160 of volume 47 of Crelle's Journal (1854), Dirichlet altered the text in several places. I have here used this later version almost throughout. K.}.}}
\end{center}

Several years ago, when the problem had arisen of transferring to the theory of complex integers the question of the number of classes of forms of the second degree corresponding to a given determinant, I had to set out the elements of the theory of forms from the beginning, since these had previously been developed only in the case of real integers. With the aid of certain considerations not previously used in this theory, and equally suited to real and complex integers, I succeeded in completing that exposition of the elements in a few pages\footnote{\emph{Recherches sur les formes quadratiques à coefficients et à indéterminées complexes}, in Crelle's Journal, vol. XXIV. Vol. I, p. 533 of this edition of G. Lejeune Dirichlet's Works. K.}. There I restricted the investigation to those properties which concern equivalence of forms, transformation, and representation of numbers by forms, since these alone were required for the question to which that paper was devoted. I did not then treat the composition of forms, an argument that the illustrious Gauss, in the fifth section of the \emph{Disquisitiones arithmeticae}, is known to have treated with the greatest generality, but by computations so lengthy that very few have been able to grasp the nature of composition, all the more because the great geometer, as he himself noted, arranged the demonstrations of the more difficult theorems synthetically in the interest of brevity, suppressing the analysis by which they had been discovered. I therefore believe I may trust that a new and wholly elementary exposition of this topic will not be unwelcome to students of analysis.

\begin{center}I.\end{center}

Before we approach the composition of forms, a few facts, either known or easily deduced from known ones, must be stated first.

Given values $\zeta,\zeta',\zeta'',\ldots$ satisfying the congruence
\[
 u^2\equiv D
\]
with respect respectively to the moduli $m,m',m'',\ldots$, we shall call them concordant with one another if a root $Z$ of the same congruence, referred to the modulus $mm'm''\ldots$, can be found such that
\[
 \zeta\equiv Z\pmod m,\qquad \zeta'\equiv Z\pmod {m'},\qquad \zeta''\equiv Z\pmod {m''},\ldots.
\]
It will suffice to consider the case in which the moduli $m,m',m'',\ldots$ are odd and prime to $D$ itself.

It is easy to see that the proposed congruences cannot be satisfied unless, with respect to each prime number dividing two or more of $m,m',m'',\ldots$, the corresponding values $\zeta,\zeta',\zeta'',\ldots$ are congruent to one another. If this condition holds, then from the given values $\zeta,\zeta',\zeta'',\ldots$ one knows the residues $\pi,\pi',\pi'',\ldots$ of $Z$ with respect to the distinct prime numbers $p,p',p'',\ldots$ dividing the product $mm'm''\ldots$. Since the residues $\pi,\pi',\pi'',\ldots$ are plainly roots of our congruence with respect respectively to the moduli $p,p',p'',\ldots$, it follows from the elementary theory of congruences that there exists a root, and in fact a unique root, $Z$ satisfying the congruence modulo $mm'm''\ldots$ and congruent to $\pi,\pi',\pi'',\ldots$ modulo $p,p',p'',\ldots$ respectively. At the same time it is immediate that
\[
 Z\equiv\zeta\pmod m,\qquad Z\equiv\zeta'\pmod {m'},\qquad Z\equiv\zeta''\pmod {m''},\ldots.
\]

Since the constant term $D$ occurring in our congruence is to retain the same value throughout what follows, we shall, for brevity, denote a given root $\zeta$ corresponding to the modulus $m$ by the notation $(m,\zeta)$. With this notation introduced, the root $Z$ to be deduced in the indicated way from the mutually concordant given roots $\zeta,\zeta',\zeta'',\ldots$, and which we shall say is composed from them, may conveniently be denoted as follows:
\[
 (m,\zeta)(m',\zeta')(m'',\zeta'')\cdots=(mm'm''\ldots,Z).
\]

We further observe that the roots $\zeta,\zeta',\zeta'',\ldots$ are always mutually concordant if $m,m',m'',\ldots$ are relatively prime to one another; and this remains true in the case that we excluded, in which $m,m',m'',\ldots$ are not prime to $2D$. Indeed, then $Z$ is already completely determined modulo $mm'm''\ldots$ by the congruences
\[
 Z\equiv\zeta\pmod m,\qquad Z\equiv\zeta'\pmod {m'},\qquad Z\equiv\zeta''\pmod {m''},\ldots,
\]
and at the same time one has
\[
 Z^2\equiv D\pmod {mm'm''\ldots}.
\]

\begin{center}II.\end{center}

For the forms of the second degree considered here,
\[
 ax^2+2bxy+cy^2,
\]
whose determinant is $D=b^2-ac$, we shall suppose the coefficients $a,b,c$ to be free of a common divisor, since the remaining cases are easily reduced to this one. It is known that forms satisfying the stated condition constitute two orders, according as at least one of the outer coefficients $a,c$ is odd or both are even. For brevity we shall exclude the latter case here, since the arguments applicable to the former apply, with the necessary changes, in the other as well.

If, in the given form, definite relatively prime values are assigned to the indeterminates $x,y$ (which will always be understood), so that
\[
 ax^2+2bxy+cy^2=m,
\]
we shall say that the number $m$ (which is always to be supposed prime to $2D$) is represented by the form. Taking integers $\xi,\eta$ such that $x\eta-y\xi=1$, it is known and easily proved that the expression
\[
 \zeta=(ax+by)\xi+(bx+cy)\eta
\]
is a root of the congruence $u^2\equiv D\pmod m$, and that the same root is always obtained in this way, however $\xi$ and $\eta$ may vary. The root $(m,\zeta)$ to which we shall say the given representation belongs can be defined in another way, more suited to our purpose. In the equation by which we have displayed $\zeta$, multiply by $y$ and put $x\eta-1$ in place of the product $y\xi$; then one sees that
\begin{equation}\tag{1}
 ax+by\equiv -y\zeta \pmod m.
\end{equation}
This congruence fully determines $\zeta$, provided that $y$ is prime to the modulus $m$. If, however, $y$ has with the modulus a greatest common divisor $\delta$ greater than unity, which plainly also divides $a$, our congruence gives only
\[
 \frac{ax+by}{\delta}\equiv -\frac{y}{\delta}\zeta\pmod {\frac m\delta},
\]
so that the residues of $\zeta$ cannot be known in this way with respect to those prime divisors of $\delta$, if there are any, which do not divide $\frac m\delta$. This inconvenience is easily remedied if one observes that the equations above imply
\[
 \zeta\equiv bx\eta,
 \qquad x\eta\equiv1\pmod\delta,
\]
and hence
\begin{equation}\tag{2}
 \zeta\equiv b\pmod\delta.
\end{equation}
Together with the preceding formula, this now fully determines the residues of $\zeta$ with respect to the individual divisors of $m$. We observe that, if $\varepsilon$ is a common divisor of $x$ and $m$, one similarly obtains
\begin{equation}\tag{3}
 \zeta\equiv -b\pmod\varepsilon.
\end{equation}

We add the following facts; although they are sufficiently well known, it is useful to put them in view here.

$1^\circ$. If two forms are equivalent, properly (as will always be understood), and the former passes into the latter by the substitution
\[
 x=\alpha x'+\beta y',\qquad y=\gamma x'+\delta y',
\]
where $\alpha\delta-\beta\gamma=1$, then clearly a number $m$ representable by one of them can also be represented by the other. It is also easy to prove that the two representations of the same number, connected by the preceding equations, belong to the same root $(m,\zeta)$. Hence the representations belonging to a given expression $(m,\zeta)$ must be referred to one whole class, and that class is unique.

$2^\circ$. Conversely, one can prove that if the same number $m$ is represented by two forms of the same determinant, and the two representations belong to the same root $(m,\zeta)$, then the forms are equivalent.

$3^\circ$. Finally, it is clear that, given any root $(m,\zeta)$, there exists a class by whose forms $m$ can be represented in such a way that the root corresponding to these representations is $(m,\zeta)$. Indeed $m$ is plainly represented by the form
\[
 \left(m,\zeta,\frac{\zeta^2-D}{m}\right),
\]
whose coefficients are free of a common divisor and in which $m$ is odd, by putting
\[
 x=1,\qquad y=0;
\]
it is evident that this representation belongs to $(m,\zeta)$.

\begin{center}III.\end{center}

With these preliminaries completed, let us approach the proposed question. Let two forms $\varphi$ and $\varphi'$ of the same determinant $D$ be given, and let $m$ and $m'$ be any two odd integers prime to $D$, representable respectively by $\varphi$ and $\varphi'$ in such a way that the roots $(m,\zeta)$ and $(m',\zeta')$ to which these representations belong are mutually concordant. Under these hypotheses, the whole matter turns on proving that the representations of $mm'$ belonging to the expression $(m,\zeta)(m',\zeta')$ will always be effected by the same form, or rather are to be referred to the same class, however $m$ and $m'$ may vary. This is the fundamental theorem in this theory.

Suppose that the given forms $(a,b,c)$ and $(a',b',c')$ have been so prepared that the expressions $(a,b)$ and $(a',b')$ are mutually concordant. This is achieved, for example, by transforming one of the forms into an equivalent one whose first coefficient has no common divisor with the first coefficient of the other. But it must be carefully noted that the analysis developed below requires nothing except that $(a,b)$ and $(a',b')$ be mutually concordant; it is not necessary that $a$ and $a'$ be relatively prime to one another, nor that they be prime to $2D$. Denote by $(aa',B)$ the expression arising from the composition of $(a,b)$ and $(a',b')$, so that
\[
 B\equiv b\pmod a,\qquad B\equiv b'\pmod {a'},\qquad D=B^2-aa'C,
\]
where $C$ is an integer. It is then clear that the forms can be changed into the following equivalent forms:
\[
 ax^2+2Bxy+a'Cy^2=\varphi,
 \qquad
 a'x'^2+2Bx'y'+aCy'^2=\varphi'.
\]
First multiplying these respectively by $a$ and $a'$, and then multiplying the resulting equations together, one obtains
\[
 (ax+By)^2-Dy^2=a\varphi,
 \qquad
 (a'x'+By')^2-Dy'^2=a'\varphi',
\]
\[
 ((ax+By)(a'x'+By')+Dyy')^2-D((ax+By)y'+(a'x'+By')y)^2=aa'\varphi\varphi'.
\]
Now, observing that $D=B^2-aa'C$, one has
\begin{equation}\tag{4}
 (ax+By)(a'x'+By')+Dyy'=aa'X+BY,
\end{equation}
where
\begin{equation}\tag{5}
 X=xx'-Cyy',\qquad Y=(ax+By)y'+(a'x'+By')y=axy'+a'x'y+2Byy'.
\end{equation}
The last equation, divided by $aa'$, assumes the form
\begin{equation}\tag{6}
 aa'X^2+2BXY+CY^2=\varphi\varphi'=\psi.
\end{equation}

After we have displayed indefinitely the product of the forms $\varphi$ and $\varphi'$ by the form $\psi$ of the same determinant $D$, suppose that both $x,y$ and $x',y'$ are assigned definite relatively prime values for which $\varphi=m$ and $\varphi'=m'$, and which satisfy the conditions indicated above. Then it remains to prove: $1^\circ$ that the representation of $mm'$ by the form $\psi$, which we see is supplied by equations (5) and (6), will be proper, that is, that $X$ and $Y$ will be free of a common divisor; and $2^\circ$ that the root to which this representation belongs, say $(mm',Z)$, will indeed be composed from the roots $(m,\zeta)$ and $(m',\zeta')$ to which the representations of $m$ and $m'$ are assumed to belong.

$1^\circ$. To show that $X$ and $Y$ have no common divisor, let us start from the supposition that a prime number $p$ divides both, and see what follows. Since $p$ must divide one of $m,m'$, let us suppose, as we may, that $m$ is a multiple of $p$. I then say that $m'$ too may be supposed divisible by $p$. Since $X$ and $Y$ are divisible by $p$, it is clear, multiplying the first of equations (5) by $-ay'$ and the second by $x'$ and adding, that an integer divisible by $p$ is obtained:
\[
 (a'x'^2+2Bx'y'+aCy'^2)y=m'y.
\]
If one were now to suppose that $m'$ is not divisible by $p$, then $y$, and consequently also $a$, would be multiples of $p$. But then, from
\[
 xx'-Cyy'=X\equiv0\pmod p,
\]
and from the fact that $x$ is prime to $y$, it follows that $x'\equiv0$, whence finally
\[
 m'=a'x'^2+2Bx'y'+aCy'^2
\]
is seen to be a multiple of $p$. Since $p$ divides both $m$ and $m'$, by equation (1) we have
\[
 ax+By\equiv -y\zeta,
 \qquad
 a'x'+By'\equiv -y'\zeta'\pmod p.
\]
Substitution of these formulas into the second of equations (5) gives the congruence
\[
 (\zeta+\zeta')yy'\equiv0\pmod p.
\]
If $\zeta+\zeta'\equiv0$, then $\zeta'\equiv-\zeta$, contrary to the hypothesis that the roots $\zeta$ and $\zeta'$ are concordant, namely that $\zeta'\equiv\zeta$. It remains to see what follows from either of the suppositions $y\equiv0$ and $y'\equiv0$, which are entirely similar. If $y$, and therefore also $x'$ by the congruence $xx'-Cyy'\equiv0$, were divisible by $p$, then by equations (2) and (3) we would have
\[
 \zeta\equiv B,
 \qquad
 \zeta'\equiv -B\pmod p,
\]
and just as above
\[
 \zeta+\zeta'\equiv0\pmod p.
\]
It is therefore established that $X$ and $Y$ are relatively prime.

$2^\circ$. We now prove that the root to which the representation of $mm'$ belongs, and which we have denoted by $(mm',Z)$, is indeed composed from $(m,\zeta)$ and $(m',\zeta')$. By symmetry it suffices to show that $Z\equiv\zeta$ with respect to any prime divisor $p$ of $m$. Since by equation (1)
\[
 ax+By\equiv -\zeta y\pmod p,
\]
equation (4) and the second of equations (5) readily give the congruences
\[
 (-\zeta(a'x'+By')+Dy')y\equiv aa'X+BY,
 \qquad
 (-\zeta y'+a'x'+By')y\equiv Y\pmod p.
\]
Multiplying the latter by $\zeta$ and adding it to the former, and taking account of the formula $\zeta^2\equiv D\pmod m$, one obtains
\[
 aa'X+BY\equiv -Y\zeta\pmod p.
\]
Comparing this congruence with
\[
 aa'X+BY\equiv -YZ\pmod p,
\]
one sees that
\[
 Z\equiv\zeta\pmod p,
\]
provided $p$ does not divide $Y$. It remains to consider the case in which $Y$ is divisible by $p$. In this case, by equation (2), $Z\equiv B\pmod p$. If now $y$ is also a multiple of $p$, then similarly $\zeta\equiv B\pmod p$, and hence
\[
 Z\equiv\zeta\pmod p.
\]
If, however, $y$ is not divisible by $p$, the congruence found above gives
\[
 a'x'+By'-\zeta y'\equiv0\pmod p,
\]
from which it is easily deduced that $p$ is a factor of the product $a'm'$. For
\[
 a'm'=(a'x'+By')^2-Dy'^2=(a'x'+By')^2-\zeta^2y'^2\equiv0\pmod p.
\]
Two cases must now be distinguished. First suppose that $a'$ is not divisible by $p$. Since in this case $m'$ is divisible by $p$, one has
\[
 a'x'+By'+\zeta'y'\equiv0\pmod p.
\]
Comparing this formula with the preceding one gives
\[
 (\zeta+\zeta')y'\equiv0\pmod p,
\]
which leads to an absurd conclusion. For since $\zeta+\zeta'$ cannot be a multiple of $p$, it would follow that $y'\equiv0\pmod p$, and hence, from the equation
\[
 m'=a'x'^2+2Bx'y'+aCy'^2,
\]
that $a'\equiv0\pmod p$, contrary to the hypothesis. Therefore this case cannot occur. In the remaining case, where $a'$ is divisible by $p$, the congruence
\[
 a'x'+By'-\zeta y'\equiv0\pmod p
\]
gives
\[
 (B-\zeta)y'\equiv0\pmod p.
\]
If $p$ does not divide $y'$, one has
\[
 \zeta\equiv B\pmod p,
\]
which agrees with the congruence $Z\equiv B\pmod p$. If, on the other hand, $p$ divides $y'$, and consequently also $m'$, then by equation (2) one has
\[
 \zeta'\equiv B\pmod p,
\]
and, as above, one concludes that $\zeta\equiv B\pmod p$, since the roots $\zeta$ and $\zeta'$ are concordant.

\end{document}
